\documentclass[11pt]{amsart}
\usepackage{amsmath,amssymb,amsthm}
\usepackage[margin=1.1in]{geometry}
\usepackage{hyperref}

\newtheorem{theorem}{Theorem}[section]
\newtheorem{lemma}[theorem]{Lemma}
\newtheorem{corollary}[theorem]{Corollary}
\newtheorem{proposition}[theorem]{Proposition}
\theoremstyle{definition}
\newtheorem{definition}[theorem]{Definition}
\theoremstyle{remark}
\newtheorem{remark}[theorem]{Remark}
\newtheorem{example}[theorem]{Example}

\newcommand{\Inf}{\mathrm{Inf}}
\newcommand{\E}{\mathbb{E}}
\newcommand{\G}{\mathsf{G}}

\title{Good coordinates in the Fourier entropy-influence conjecture}
\author{Yitzchak Shmalo}
\date{}

\begin{document}

\begin{abstract}
For a Boolean function $f$ let $E_i(f)=H(\mathbf 1[i\in U]\mid U\setminus\{i\})$, where $U$ is drawn from the spectral
distribution $\hat f^2$. The Fourier entropy-influence conjecture holds if and only if the quantity
$\G=\sup_f\min_i E_i(f)/\Inf_i(f)$ is finite, and then its constant lies between $\G$ and $\G+2/\ln 2$. We study $\G$
directly. A composition identity computes $E_i$ and $\Inf_i$ for iterated block compositions $g(h,\dots,h)$ exactly
from the biased spectrum of $g$, and towers of transitive monotone gates give $\G_{\mathrm{mon}}\ge 4.50642846$,
where $\G_{\mathrm{mon}}$ is the supremum over monotone functions. On the other side we show that every width-$w$ DNF
has a coordinate with $E_i/\Inf_i\le 2+\log_2 e+\log_2(2^{w+1}-1)$, with no monotonicity or read restriction, and
that every relevant leaf of a read-once AND/OR formula has $E_i/\Inf_i<\pi/\ln 2$, which is sharp. Finally we show
that the coordinate of largest influence is not a valid selection rule, and that all but a vanishing fraction of the
coordinates of a monotone function can have $E_i/\Inf_i$ bounded below by any constant. Nor can the ratios be
averaged with weights $\Inf_i^{a+1}$ for any $a\ne0$.
\end{abstract}

\maketitle

\section{Introduction}

Let $f:\{-1,1\}^n\to\{-1,1\}$ have Fourier expansion $f=\sum_S\hat f(S)\chi_S$, and let $U\subseteq[n]$ be drawn with
$\Pr[U=S]=\hat f(S)^2$. The spectral entropy of $f$ is $H(f)=H(U)$ and the influence of $i$ is
$\Inf_i(f)=\Pr[i\in U]$. The Fourier entropy-influence (FEI) conjecture of Friedgut and Kalai~\cite{FK} asserts
that $H(f)\le C\sum_i\Inf_i(f)$ for a universal constant $C$; let $C^\ast$ be the least such constant, $C^\ast=\infty$
if there is none. The conjecture is known for several classes~\cite{OWZ,CKLS,KMS,WWW14,Pot26} and with a
$\log$-type loss in general~\cite{Han,KMS}; Kelman, Kindler, Lifshitz, Minzer and Safra~\cite{KKLMS} prove a weaker
form, that the spectrum of a Boolean function of constant variance is concentrated on $2^{O(\mathrm I(f)\log\mathrm
I(f))}$ characters. The largest known lower bound is $C^\ast\ge 6.4548$~\cite{Hod}.

Gonz\'alez, MacManus and Pereyra~\cite{GMP} reduce the conjecture to a statement about one coordinate: it is enough
that every Boolean function has a coordinate $i$ at which the entropy drop
$H(f)-\tfrac12H(f_{i\to1})-\tfrac12H(f_{i\to-1})$ is at most $K\Inf_i(f)$. They also observe that the choice of
coordinate matters: for their $\delta$-tribes functions a stronger form of the inequality fails at one coordinate
and holds, with constant $38$, at another~\cite[Section~5]{GMP}. In~\cite{Shm2} it is shown that the reduction is
lossless up to an additive $2/\ln 2$, in the following form. Write
\[
E_i(f)=H\bigl(\mathbf 1[i\in U]\bigm| U\setminus\{i\}\bigr),\qquad
\G(f)=\min_{i:\,\Inf_i(f)>0}\frac{E_i(f)}{\Inf_i(f)},\qquad \G=\sup_f\G(f),
\]
the supremum over all nonconstant Boolean functions on any number of variables. Then
\[
\G\le C^\ast\le\G+\frac{2}{\ln 2}.
\]
So the conjecture is equivalent to $\G<\infty$, and any proof of it that proceeds by choosing one good coordinate
at a time needs a selection rule whose ratio $E_i/\Inf_i$ is bounded. This paper studies $\G$ itself: how large it
must be, which classes have a coordinate with bounded ratio, and which natural selection rules fail.

\subsection*{Results}
Section~\ref{sec:comp} gives an exact formula for $E_i$ and $\Inf_i$ under block composition $f=g(h,\dots,h)$
(Theorem~\ref{thm:comp}). The formula carries four numbers per leaf through each level of a composition, so towers
of depth ten or more on $10^{12}$ and more variables are computed in milliseconds and to any precision.

Section~\ref{sec:towers} uses it to bound $\G_{\mathrm{mon}}$, the supremum of $\G(f)$ over monotone $f$, from below.
When every gate is transitive on its inputs, every leaf of the tower has the same ratio, so the ratio of one leaf is
$\G(f)$. Towers of the perfect matching function of $K_{3,3}$ and its dual already give $\G_{\mathrm{mon}}>4.505$;
mixed words of six gate types give
\[
\G\ge\G_{\mathrm{mon}}\ge 4.5064284654452331779\ldots
\]
This is below the bound $C^\ast\ge 6.4548$, as it must be: towers are built to make one coordinate's ratio large,
and the spectral entropy of these functions is far from extremal. The point of the bound is the other one: any
proof of FEI through a good coordinate must allow ratio at least $4.506$ already for monotone functions.

Section~\ref{sec:classes} bounds $\G(f)$ from above for classes. Every coordinate satisfies
$E_i/\Inf_i\le 2+\log_2 e+\log_2(1+\mathrm{Var}(f)/\Inf_i)$, and from this every width-$w$ DNF or CNF has a
coordinate with ratio at most $2+\log_2 e+\log_2(2^{w+1}-1)$. For read-once AND/OR formulas every relevant leaf
has ratio below $\pi/\ln 2$, and the constant is attained in the limit.

Section~\ref{sec:neg} gives three negative results. The coordinate of largest influence can have unbounded ratio,
even for monotone functions. In a monotone function built from a middle-layer slice all but $k$ of
$\binom{k}{\lfloor k/2\rfloor}+k$ coordinates can have ratio above any prescribed constant. And the weighted
inequalities $\sum_i\Inf_i^{a}E_i\le C\sum_i\Inf_i^{a+1}$, which would give a selection rule by averaging, fail for
every $a\ne0$. The case $a=0$ is the conjecture itself.

All numerical values in this paper are reproduced by the scripts in the accompanying repository, in double
precision and with $40$-digit arithmetic.

\section{Composition}\label{sec:comp}

\subsection{Biased expansion of the gate}
Let $g:\{-1,1\}^k\to\{-1,1\}$ and $\mu\in(-1,1)$. Put $\sigma=\sqrt{1-\mu^2}$ and $\phi_\mu(y)=(y-\mu)/\sigma$. If
$y_1,\dots,y_k$ are independent with mean $\mu$, the products $\prod_{l\in T}\phi_\mu(y_l)$ are orthonormal, and
\[
g(y)=\sum_{T\subseteq[k]}\hat g_\mu(T)\prod_{l\in T}\phi_\mu(y_l),\qquad
\hat g_\mu(T)=\E\Bigl[g(y)\prod_{l\in T}\phi_\mu(y_l)\Bigr],\qquad \sum_T\hat g_\mu(T)^2=\E[g^2]=1.
\]
In particular $\hat g_\mu(\emptyset)=\E[g(y)]$.

\subsection{The spectral sample of a composition}
Let $h:\{-1,1\}^m\to\{-1,1\}$ have mean $\mu=\hat h(\emptyset)\in(-1,1)$, and let
$f(x^1,\dots,x^k)=g(h(x^1),\dots,h(x^k))$ on $k$ disjoint blocks of $m$ variables. Since
$\phi_\mu(h)=\sigma^{-1}\sum_{A\ne\emptyset}\hat h(A)\chi_A$, expanding $g$ gives
\[
\hat f\Bigl(\bigcup_{l\in T}A_l\Bigr)=\hat g_\mu(T)\prod_{l\in T}\frac{\hat h(A_l)}{\sigma}
\qquad(A_l\ne\emptyset\text{ a subset of block }l),
\]
and every Fourier coefficient of $f$ has exactly one such representation. Since
$\sum_{A\ne\emptyset}\hat h(A)^2=1-\mu^2=\sigma^2$, the spectral sample of $f$ is the following two-stage draw: choose
$T\subseteq[k]$ with probability $\hat g_\mu(T)^2$, then, independently for each $l\in T$, choose a nonempty
$A_l$ in block $l$ with probability $\hat h(A_l)^2/\sigma^2$, and set $U=\bigcup_{l\in T}A_l$.

\subsection{The leaf state}
Fix a coordinate $i$ of $h$. Write $m(a,b)=(a+b)\,h_2\bigl(b/(a+b)\bigr)$ for $a,b\ge 0$, where
$h_2(p)=-p\log_2p-(1-p)\log_2(1-p)$ and $m(0,0)=0$; this is the unnormalised entropy of a two-point mass. Grouping
the spectral sets of $h$ according to $B=U\setminus\{i\}$,
\begin{equation}\label{eq:Esplit}
E_i(h)=\sum_{B\not\ni i}m\bigl(\hat h(B)^2,\hat h(B\cup\{i\})^2\bigr)=e(h)+m\bigl(\mu^2,\hat h(\{i\})^2\bigr),
\end{equation}
where $e(h)$ is the sum over $B\ne\emptyset$. We call $(\mu,\,\hat h(\{i\})^2,\,e(h),\,\Inf_i(h))$ the state of the
leaf $i$ in $h$.

\begin{theorem}\label{thm:comp}
Let $f=g(h,\dots,h)$ as above, let $i$ be a coordinate of $h$ placed in block $j$, and let $(\mu,P,e,I)$ be its
state in $h$. Write $\hat g$ for $\hat g_\mu$ and $W_j=\sum_{T\ni j}\hat g(T)^2$, the $\mu$-biased influence of $j$
on $g$. Then the state $(\mu',P',e',I')$ of $i$ in $f$ is
\begin{align*}
\mu'&=\hat g(\emptyset), & P'&=\hat g(\{j\})^2\,\frac{P}{\sigma^2},\\
I'&=W_j\,\frac{I}{\sigma^2}, & e'&=W_j\,\frac{e}{\sigma^2}
+\sum_{\substack{R\ne\emptyset\\ j\notin R}} m\Bigl(\hat g(R)^2,\ \hat g(R\cup\{j\})^2\,\frac{P}{\sigma^2}\Bigr),
\end{align*}
and $E_i(f)=e'+m(\mu'^2,P')$.
\end{theorem}

\begin{proof}
$\mu'=\E f=\hat g(\emptyset)$, and $P'=\hat f(\{i\})^2$ is the case $T=\{j\}$, $A_j=\{i\}$ of the product formula.
For $I'$: $i\in U$ exactly when $j\in T$ and $i\in A_j$, and these events are independent with probabilities $W_j$
and $\sum_{A\ni i}\hat h(A)^2/\sigma^2=I/\sigma^2$.

For $E_i(f)$ we condition on $U\setminus\{i\}$, which determines $R=T\setminus\{j\}$, the sets $A_l$ for $l\ne j$,
and $B=A_j\setminus\{i\}$ when $j\in T$. Whether $i\in U$ depends only on what happens in block $j$.

If $B\ne\emptyset$ then $j\in T$, and given $B$ the event $i\in A_j$ has the conditional law it has in the
normalised nonempty spectrum of $h$; the blocks $l\ne j$ are independent of it. Summing over $T\ni j$ and over
nonempty $B$ gives $W_j\,e/\sigma^2$.

If $B=\emptyset$ then either $j\notin T$, so $T=R$ and $i\notin U$, with probability $\hat g(R)^2$, or $T=R\cup\{j\}$
and $A_j=\{i\}$, with probability $\hat g(R\cup\{j\})^2P/\sigma^2$. The outcomes in the other blocks factor out of
both alternatives with the same weight, so this part contributes
$\sum_{R\not\ni j}m(\hat g(R)^2,\hat g(R\cup\{j\})^2P/\sigma^2)$. The term $R=\emptyset$ is the one in which
$U\setminus\{i\}=\emptyset$; it equals $m(\mu'^2,P')$, and the remaining terms together with $W_je/\sigma^2$ form
$e'$, the part of $E_i(f)$ over nonempty $U\setminus\{i\}$.
\end{proof}

A single variable $x_i$ has state $(0,1,0,1)$. Iterating Theorem~\ref{thm:comp} computes the state of any leaf of an
iterated composition $g_d(\dots g_2(g_1(x,\dots),\dots)\dots)$, and the cost of a level is one biased transform of
the gate, $O(k2^k)$, independent of the number of variables below it. We checked the formula against brute-force
Fourier transforms on all $444$ leaf positions of depth-two compositions of gates on three and four variables, with
largest discrepancy $3.3\times10^{-16}$.

\begin{remark}
If $g$ is transitive on its inputs, so is the composition on its leaves, and every leaf has the same state. For such
towers the ratio of one leaf is $\G(f)$ and also $\max_iE_i/\Inf_i$; a lower bound on $\G$ then needs no search over
coordinates. Without transitivity the formula still applies to each leaf separately, with $W_j$ and $\hat g(\{j\})$
taken at that leaf's position.
\end{remark}

\begin{remark}
Keller, Mossel and Schlank~\cite[Claim~4.1]{KMS} bound $H(f)$ above by $\sum_i h_2(\Inf_i)$ using subadditivity of
entropy over the indicators $\mathbf 1[i\in U]$. The quantity $E_i$ is the conditional counterpart: it measures how
much of the uncertainty of $\mathbf 1[i\in U]$ remains once the rest of $U$ is known.
\end{remark}

\section{Towers of monotone gates}\label{sec:towers}

\subsection{Gates and convention}
We write a gate as a monotone function of $z\in\{0,1\}^k$ and pass to $\pm1$ values by $x=-1$ when $z=1$ and
output $-1$ when the gate is true. The gates used are:
\begin{itemize}
\item $M$, the perfect-matching function of $K_{3,3}$: $z$ indexes the nine edges and $M(z)=1$ when the chosen edges
contain a perfect matching; $D(z)=1-M(\mathbf 1-z)$ is its dual;
\item $P$, the perfect-matching function of $K_{4,4}$ on sixteen edges, and its dual $P^\ast$;
\item the tribes functions $T(w,s)$, the OR of $s$ disjoint ANDs of $w$ variables each, and their duals
$T(w,s)^\ast$; we also write $Q_2=T(2,2)^\ast$, and $R^\ast$ for the dual of the function on the six edges of $K_4$ that is
true when two opposite edges are both present (a copy of $T(2,3)$).
\end{itemize}
Each gate is monotone and its automorphism group is transitive on its inputs. A word $g_1g_2\cdots g_d$ denotes the
iterated composition in which $g_1$ is applied to disjoint blocks of variables, $g_2$ to disjoint blocks of copies
of the result, and so on; a word is read from the bottom gate to the top gate. Compositions of monotone functions
are monotone, and since every gate is transitive the composition is transitive on its leaves, so
$\G(f)=E_i(f)/\Inf_i(f)$ for any leaf $i$. That ratio is computed by iterating Theorem~\ref{thm:comp} from the
state $(0,1,0,1)$ of a single variable.

\begin{proposition}\label{prop:towers}
For the words in Table~\ref{tab:towers} the iterated compositions are monotone Boolean functions with
$\G(f)$ equal to the value shown, to the number of digits shown. In particular
\[
\G_{\mathrm{mon}}\ge 4.50642846544523317790928639\ldots
\]
\end{proposition}

\begin{table}[h]
\centering
\begin{tabular}{lrl}
word (bottom to top) & variables & $\G(f)$\\\hline
$M\,D\,D\,D\,D\,T(2,5)^\ast$ & $5.9\times10^{5}$ & $4.50585006207690030206437435306$\\
$M\,D\,D\,M\,M\,D\,D$ & $4.8\times10^{6}$ & $4.50504996749762792695710133417$\\
$M\,D\,D\,D\,D\,D\,M\,M$ & $4.3\times10^{7}$ & $4.50637987946771061063535681398$\\
$D\,M^{12}$ & $2.5\times10^{12}$ & $4.50642830552543541108624856879$\\
$M\,D^4\,M\,P\,P\,P\,Q_2\,R^\ast\,T(4,3)^\ast\,T(2,6)\,T(3,5)$ & $1.1\times10^{14}$ & $4.50640064551556843550717367045$\\
$D\,M^{8}\,P^\ast\,M\,M\,T(3,5)$ & $7.5\times10^{12}$ & $4.50642846293912877531310941280$\\
$M\,D^{8}\,P\,D\,D\,T(4,3)^\ast\,T(3,5)\,T(3,5)$ & $1.4\times10^{15}$ & $4.50642846544523317790928639245$
\end{tabular}
\caption{Lower bounds for $\G_{\mathrm{mon}}$ from towers. Values to 30 digits are computed with 40-digit arithmetic.}
\label{tab:towers}
\end{table}

\begin{proof}
Each value is the output of Theorem~\ref{thm:comp} iterated along the word, with the biased coefficients of each gate
computed from its truth table. The computation is exact up to rounding, and was carried out with 40 significant
digits; the digits shown are stable under a change to 60. For $D\,M^{12}$ and the last two rows the iteration was
also run in 60-digit interval arithmetic with outward rounding; each enclosure has width below $10^{-48}$ and
contains the value shown. Theorem~\ref{thm:comp} itself was checked against
brute-force Fourier transforms as described at the end of Section~\ref{sec:comp}.
\end{proof}

\begin{remark}
Double precision is not enough to rank long towers. For the last word of Table~\ref{tab:towers} the double-precision
iteration returns $4.506428469267$, an error of $3.8\times10^{-9}$, which is larger than the gap of
$2.5\times10^{-9}$ between the last two rows. The ordering of the table is that of the 40-digit values.
\end{remark}

\begin{remark}[Where the value comes from]
For $M$ the mean map $\mu\mapsto\hat M_\mu(\emptyset)$ has a fixed point at $\mu^\ast=-0.0349185099940\ldots$, where
its derivative is $2.0077$, so the fixed point is repelling. The bottom gate $D$ produces the mean
$-0.03515625$, close to $\mu^\ast$, and a run of $M$ gates then keeps the mean near $\mu^\ast$ for several levels
before it escapes; the top gates spend what is left. If the mean could be held at $\mu^\ast$ for ever, starting from
the state produced by $D$, the ratio would tend to $4.5064554626637$. This is a heuristic target rather than a
bound: no finite tower holds the mean at a repelling fixed point, the limit depends on the starting state
($4.5064428$ from the state after $D\,M$, $4.5064353$ after $D\,M\,M$), and we do not know whether
$\G_{\mathrm{mon}}$ exceeds it. The same heuristic for other gates gives larger targets, $4.5168$ for $T(2,2)^\ast$
pinned at its fixed mean $0.236068$ and $4.5107$ for $T(4,3)^\ast$, but towers that bring the mean close to these
fixed points carry a smaller leaf ratio, and none of depth at most four below the pinned gate exceeds $4.504$.
\end{remark}

\section{Classes with a good coordinate}\label{sec:classes}

Throughout this section $f$ is Boolean, $i$ is a coordinate with $\delta=\Inf_i(f)>0$, and for $T\not\ni i$ we put
$p_T=\hat f(T)^2$ and $q_T=\hat f(T\cup\{i\})^2$. Conditioning on $U\setminus\{i\}=T$ gives
\begin{equation}\label{eq:erasure}
E_i(f)=\sum_{T\not\ni i}m(p_T,q_T),\qquad \Inf_i(f)=\sum_{T\not\ni i}q_T .
\end{equation}
The function $m$ is concave, increasing in each argument and homogeneous of degree one, hence superadditive:
$\sum_T m(a_T,b_T)\le m(\sum_Ta_T,\sum_Tb_T)$. Write $g=\E_if$ and $h=\mathrm D_if$, so that $f=g+x_ih$,
$p_T=\hat g(T)^2$ and $q_T=\hat h(T)^2$; since $h$ takes values in $\{-1,0,1\}$, $\sum_Tq_T=\E h^2=\delta$.

\begin{lemma}\label{lem:var}
For every Boolean $f$ and every coordinate $i$ with $\delta=\Inf_i(f)>0$,
\[
\frac{E_i(f)}{\delta}\le 2+\log_2e+\log_2\Bigl(1+\frac{\mathrm{Var}(f)}{\delta}\Bigr).
\]
\end{lemma}

\begin{proof}
For $t\in[0,1]$ one has $h_2(t)\le 2\sqrt{t(1-t)}$, so $m(a,b)\le 2\sqrt{ab}$. The term $T=\emptyset$ is therefore at
most $2\sqrt{q_\emptyset}=2|\E h|\le 2\delta$. For the others, $\sum_{T\ne\emptyset}p_T=\mathrm{Var}(g)\le\mathrm{Var}(f)=V$
and $\sum_{T\neq\emptyset}q_T\le\delta$, so by superadditivity and monotonicity
\[
\sum_{T\ne\emptyset}m(p_T,q_T)\le m(V,\delta)=V\log_2\Bigl(1+\frac{\delta}{V}\Bigr)+\delta\log_2\Bigl(1+\frac V\delta\Bigr)
\le\frac{\delta}{\ln2}+\delta\log_2\Bigl(1+\frac V\delta\Bigr).\qedhere
\]
\end{proof}

So a coordinate whose influence is a fixed fraction of the variance is good. Width supplies one.

\begin{theorem}\label{thm:width}
Let $f$ be a nonconstant function computed by a DNF in which every term has at most $w$ literals. Then $f$ has a
coordinate with
\[
\frac{E_i(f)}{\Inf_i(f)}\le 2+\log_2e+\log_2\bigl(2^{w+1}-1\bigr)<w+3+\log_2e .
\]
The same holds for CNFs of width $w$. No monotonicity and no bound on the number of occurrences of a variable is
assumed.
\end{theorem}

\begin{proof}
Since $f$ is not constant, some term is satisfiable and nonempty; let its literals be $\ell_1,\dots,\ell_k$, $k\le w$,
on coordinates $j_1,\dots,j_k$. From every input $x$ with $f(x)$ false, set $\ell_1,\dots,\ell_k$ true one at a time.
The last input satisfies the term, so some step crosses an edge of the cube along which $f$ changes. An edge in
direction $j_t$ is reached at step $t$ from at most $2^{t-1}$ starting inputs, those that differ from its lower end
only in $j_1,\dots,j_{t-1}$. The number of such edges in direction $j$ is $2^{n-1}\Inf_j(f)$. Hence
\[
\Pr[f\text{ false}]\le\frac12\sum_{t=1}^k2^{t-1}\Inf_{j_t}(f)\le\frac{2^k-1}{2}\max_t\Inf_{j_t}(f).
\]
Take $i$ the maximiser. With $V=\mathrm{Var}(f)=4\Pr[f\text{ false}]\Pr[f\text{ true}]$ this gives
$V/\Inf_i(f)\le2(2^k-1)$, and Lemma~\ref{lem:var} gives the bound. For a CNF apply this to $\neg f$, which has the
same $E_i$ and $\Inf_i$.
\end{proof}

Width is preserved by restriction, so the induction of~\cite{Shm2} turns Theorem~\ref{thm:width} into
$H(f)\le(w+3+3\log_2e)\,\mathrm I(f)$ for width-$w$ DNFs and CNFs. The argument has the same shape as the
coordinate choice of~\cite[Section~5.2]{GMP}: a variable inside a short term sees a positive fraction of the
variance. For comparison, Han~\cite{Han} proves $H(f)=O\bigl(\mathrm I(f)+\sum_i\Inf_i(f)\log(1/\Inf_i(f))\bigr)$
for every $f$, and Arunachalam, Chakraborty, Kouck\'y, Saurabh and de Wolf~\cite{ACKSW} bound $H(f)$ by twice the
average unambiguous parity-certificate complexity.

\subsection{Read-once formulas}
For read-once formulas a much stronger statement holds: every relevant leaf is good, with a sharp constant.

\begin{theorem}\label{thm:readonce}
Let $f$ be a read-once formula over AND and OR with literals at the leaves. Every leaf $i$ with $\Inf_i(f)>0$
satisfies
\[
\frac{E_i(f)}{\Inf_i(f)}<\frac{\pi}{\ln 2}=4.5323601418\ldots
\]
and the constant $\pi/\ln 2$ cannot be lowered: there are monotone read-once formulas and leaves with ratio
arbitrarily close to it.
\end{theorem}

Negating an input or the output changes neither $E_i$ nor $\Inf_i$, so we may assume $f$ monotone. Let the gates on
the path from leaf $i$ to the root be numbered $1,\dots,L$ upwards. At gate $l$ the other children form a read-once
formula on a block of fresh variables; let $Z_l\in\{0,1\}$ be the event that they let the path child through
(all true at an AND, all false at an OR), $p_l=\E Z_l$, and $c_l\in\{-1,1\}$ the output when they block it. We may
assume $0<p_l<1$: a gate with $p_l=1$ can be deleted and one with $p_l=0$ makes $i$ irrelevant. Then
$f_0=x_i$, $f_l=c_l(1-Z_l)+f_{l-1}Z_l$ and $f=f_L$. Put $P_0=1$, $P_l=p_1\cdots p_l$ and $\delta=P_L=\Inf_i(f)$.

\begin{lemma}\label{lem:path}
Let $\Phi(u)=m(u^2,1)$. With
$R_r=\sum_{l<r}c_l(1-p_l)/P_l-c_r/P_{r-1}$ and $R_\infty=\sum_{l\le L}c_l(1-p_l)/P_l$,
\[
\frac{E_i(f)}{\delta}=\sum_{r=1}^L(P_{r-1}-P_r)\,\Phi(R_r)+\delta\,\Phi(R_\infty).
\]
\end{lemma}

\begin{proof}
Unrolling the recursion, $h=\mathrm D_if=\prod_lZ_l$ and $g=\E_if=\sum_lc_l(1-Z_l)\prod_{j>l}Z_j$. A frequency
$S\not\ni i$ splits into block parts $S_l$, and $\hat h(S)=\prod_l\hat Z_l(S_l)$ with $\hat Z_l(\emptyset)=p_l$. If
$r$ is the first $l$ with $S_l\ne\emptyset$, only the terms $l\le r$ of $g$ contribute, and
$\hat g(S)=R_r\hat h(S)$; if all $S_l$ are empty, $\hat g(S)=R_\infty\hat h(S)$. By homogeneity
$m(p_S,q_S)=q_S\Phi(\hat g(S)/\hat h(S))$ when $\hat h(S)\ne0$, and $m(p_S,0)=0$. Under the law $q_S/\delta$ the blocks
are independent, block $l$ is empty with probability $p_l$, so the first nonempty block is $r$ with probability
$P_{r-1}-P_r$ and none is nonempty with probability $\delta$. Now apply~\eqref{eq:erasure}.
\end{proof}

\begin{proof}[Proof of Theorem~\ref{thm:readonce}]
Let $a_l\le b_l$ be the probabilities that $f_l$ is true when $x_i$ is false, respectively true, and
$d_l=b_l-a_l=P_l$. Put $\alpha=a/d$, $\beta=(1-b)/d$, so $s=1+\alpha+\beta=1/d$. An AND gate keeps $\alpha$ and raises
$\beta$; an OR gate keeps $\beta$ and raises $\alpha$. In terms of the values $(\alpha,\beta)$ before gate $r$,
$\sum_{l<r}c_l(1-p_l)/P_l=\beta-\alpha$ (with $c=-1$ for true), so $|R_r|=1+2\alpha$ at an AND and $|R_r|=1+2\beta$ at an
OR, and $|R_\infty|=|\beta-\alpha|$ at the end. With $\psi(t)=\Phi(2t+1)$, the charge of an AND step is
\[
(P_{r-1}-P_r)\,\psi(\alpha)=\int_{\beta}^{\beta'}\frac{\psi(\alpha)}{(1+\alpha+\beta)^2}\,d\beta ,
\]
and symmetrically at an OR. The function $\Phi$ is increasing on $[0,\infty)$, and for $u\ge1$
\[
\Phi''(u)\ln2=2\ln(1+u^{-2})-\frac{4}{1+u^2}<0,
\]
since $\ln(1+x)<2x/(1+x)$ for $0<x\le1$. So $\psi$ is increasing and strictly concave on $[0,\infty)$. For $\alpha\ge\beta$ let
\[
W(\alpha)=\int_\alpha^\infty\frac{2\psi(t)}{(1+2t)^2}\,dt,\qquad
V(\alpha,\beta)=\psi(\alpha)\Bigl[\frac1s-\frac1{1+2\alpha}\Bigr]+W(\alpha),
\]
and $V(\alpha,\beta)=V(\beta,\alpha)$ for $\alpha<\beta$; $V$ is continuous. For $\alpha\ge\beta$, $-V_\beta=\psi(\alpha)/s^2$, and
\[
-V_\alpha-\frac{\psi(\beta)}{s^2}=\frac{\psi(\alpha)-\psi(\beta)}{s^2}-\frac{\psi'(\alpha)(\alpha-\beta)}{s(1+2\alpha)}
\ge \psi'(\alpha)(\alpha-\beta)\Bigl[\frac1{s^2}-\frac1{s(1+2\alpha)}\Bigr]\ge0
\]
by concavity and $s\le1+2\alpha$. Hence along every step the charge is at most the decrease of $V$, on either side of
the diagonal. At the end, $W(\alpha)>\psi(\alpha)/(1+2\alpha)$ because $\psi$ is increasing, so for $\alpha\ge\beta$
\[
V(\alpha,\beta)>\frac{\psi(\alpha)}{s}\ge\frac{\Phi(\alpha-\beta)}{s}=\delta\,\Phi(R_\infty).
\]
Lemma~\ref{lem:path} and the start $(\alpha,\beta)=(0,0)$ give $E_i/\delta<V(0,0)=W(0)=\int_1^\infty\Phi(u)u^{-2}\,du$.
Since $\Phi'(u)=2u\ln(1+u^{-2})/\ln2$ and $\Phi(1)=2$, integration by parts and the antiderivative
$u\ln(1+u^{-2})+2\arctan u$ give
\[
\int_1^\infty\frac{\Phi(u)}{u^2}\,du=2+\frac{2}{\ln 2}\int_1^\infty\ln(1+u^{-2})\,du=2+\frac2{\ln2}\Bigl(\frac\pi2-\ln2\Bigr)
=\frac{\pi}{\ln2}.
\]

For sharpness take $q=1-2^{-m}$ and $K=m2^m$ gates alternating AND, OR, AND, \dots; at an AND the sibling is an OR of
$m$ fresh variables and at an OR an AND of $m$ fresh variables, so every $p_l=q$. Then $\alpha$ and $\beta$ differ by
$O(1-q)$ relative to $s$, every step multiplies $s$ by $1/q$, and $s$ runs from $1$ to $q^{-K}\approx e^m$. The
charges are Riemann sums of $\int_1^{q^{-K}}\Phi(u)u^{-2}\,du$ with mesh tending to zero, and the terminal term is
nonnegative, so the ratio tends to $\pi/\ln2$ as $m\to\infty$. These formulas have $1+m^22^m$ variables.
\end{proof}

The values for $m=2,4,\dots,14$ are $3.633$, $4.326$, $4.492$, $4.525$, $4.531$, $4.5322$, $4.5323$, and the
identity of Lemma~\ref{lem:path} agrees with brute-force transforms of random guarded paths to $3\times10^{-15}$.
Read-once formulas are known to satisfy FEI~\cite{OWZ,CKLS,OT13}, and for the whole one-coordinate step
$H(f)-\frac12H(f_{i\to1})-\frac12H(f_{i\to-1})$ every leaf of a read-once formula has ratio below $5+\ln4$ in
nats~\cite{Shm1}; the
content of Theorem~\ref{thm:readonce} is that the erasure ratio is bounded at every leaf, however small its influence,
with the exact constant.

\section{Selection rules that fail}\label{sec:neg}

All three constructions below use one device, a single point added to a monotone function, which is also the source of
the example in~\cite[Proposition~5.1]{GMP}. Let $k=2m+1$, write $\mathrm{Maj}$ for majority on $k$ variables with value
$-1$ when true, and let $\hat{\mathrm{Maj}}$ denote its Fourier coefficients. For odd $|T|=s=2j+1$,
\[
|\hat{\mathrm{Maj}}(T)|=\frac{\binom{2m}{m}}{4^m}\cdot\frac{\binom mj}{\binom{2m}{2j}}
\]
\cite[Section~5.3]{ODo}. From $\binom{2m}m\ge4^m/(2m+1)$, $\binom mj\ge2^{mh_2(j/m)}/(m+1)$ and
$\binom{2m}{2j}\le2^{2mh_2(j/m)}$,
\begin{equation}\label{eq:majlow}
2^k|\hat{\mathrm{Maj}}(T)|\ge\frac{2^{m+1}}{(2m+1)(m+1)}\qquad(|T|\text{ odd}).
\end{equation}

\subsection{The coordinate of largest influence}

\begin{proposition}\label{prop:maxinf}
For every $C$ there is a monotone Boolean function whose coordinate of largest influence is unique and has
$E_i/\Inf_i>C$.
\end{proposition}

\begin{proof}
Fix an $m$-subset $A$ of $[k]$ and let $F(x,y)=\mathrm{Maj}(y)\vee(x\wedge\bigwedge_{j\in A}y_j)$. The only $y$ at which
$x$ matters is $y=1_A$, so $\Inf_x(F)=2^{-k}$, $\mathrm D_xF=\pm\mathbf 1[y=1_A]$ and $q_T=4^{-k}$ for all $T\subseteq[k]$.
Moreover $\E_xF=\mathrm{Maj}-\mathrm{Maj}(1_A)\mathbf 1[y=1_A]$ differs from $\mathrm{Maj}$ in each coefficient by $2^{-k}$.
Since $m(a,b)\ge b\log_2(1+a/b)$, \eqref{eq:erasure} and~\eqref{eq:majlow} give, for $m$ large,
\[
\frac{E_x(F)}{\Inf_x(F)}\ge2^{-k}\sum_{|T|\text{ odd}}\log_2\bigl(1+(2^k\hat F(T))^2\bigr)
\ge\log_2\Bigl(\frac{2^{m+1}}{(2m+1)(m+1)}-1\Bigr)=\frac k2-O(\log k).
\]
Now replace each $y_j$ by a majority of $b$ fresh variables, $b$ odd. Because the inner function is balanced, the
spectral sample of the composition is that of $F$ with each block refined independently by $\hat{\mathrm{Maj}_b}^2$,
which has total mass one on nonempty sets; by homogeneity of $m$ this leaves $E_x$ and $\Inf_x$ unchanged, while each
new leaf has influence $\Inf_{y_j}(F)\binom{b-1}{(b-1)/2}2^{1-b}\to0$. For $b$ large $x$ is the unique coordinate of
largest influence, and the composition is monotone.
\end{proof}

For $k=3,5,\dots,17$ the ratio of $x$ is $3.52$, $4.24$, $4.99$, $5.78$, $6.62$, $7.48$, $8.37$, $9.27$.

\subsection{Most coordinates hidden}

\begin{proposition}\label{prop:middle}
For $k=2m+1$ let
\[
F(y,z)=[\,|y|\ge m+1\,]\ \vee\bigvee_{|A|=m}\Bigl(z_A\wedge\bigwedge_{j\in A}y_j\Bigr),
\]
a monotone function of $k+\binom km$ variables. Each of the $\binom km$ coordinates $z_A$ has influence $2^{-k}$ and
\[
\frac{E_{z_A}(F)}{\Inf_{z_A}(F)}=2^{-k}\sum_{T\subseteq[k]}\Phi\bigl(2^k\hat a(T)\bigr)\ge\frac k2-O(\log k),
\]
where $a=\mathrm{Maj}-\mathbf 1[|y|=m]$ and $\hat a(T)=\bigl(1-\tfrac{s}{k+1}\bigr)\hat{\mathrm{Maj}}(T)$ for odd $|T|=s$.
\end{proposition}

\begin{proof}
On the layer $|y|=m$ the output is $z_y$, below it false and above it true. So $F=a(y)\pm\sum_A z_A\mathbf 1[y=1_A]$
with no term containing two hidden variables, and for $z_A$ one has $q_T=4^{-k}$ and $p_T=\hat a(T)^2$ on
$T\subseteq[k]$, while frequencies containing another $z_B$ have $q=0$. This gives the identity. For the coefficients,
let $\mathcal L$ be the Laplacian, $\mathcal L\chi_T=|T|\chi_T$. Majority changes only between the layers $m$ and $m+1$,
across $m+1$ edges from each point, so $\mathcal L\,\mathrm{Maj}=(m+1)(L(y)-L(\bar y))$ with $L=\mathbf 1[|y|=m]$. At odd
$T$ this reads $s\,\hat{\mathrm{Maj}}(T)=(k+1)\hat L(T)$. The lower bound follows from~\eqref{eq:majlow}, $1-s/(k+1)\ge1/(k+1)$
and $\Phi(u)\ge\log_2(1+u^2)$, as before.
\end{proof}

So all but $k$ of the $k+\binom km$ coordinates have unbounded ratio. The $k$ base coordinates stay good. With
$r=\binom km$ and $\delta=2^{-k}$, a base coordinate $y_1$ has influence $r\delta$. A frequency containing one hidden
variable $z_B$ has coefficient $\pm2^{-k}\chi_T(1_B)$ at $T\cup\{z_B\}$, $T\subseteq[k]$, so the pairs $T$,
$T\cup\{1\}$ have equal squares $4^{-k}$, and since $m(c,c)=2c$ these frequencies contribute exactly $r\delta$ to
$E_{y_1}$. The remaining frequencies give $E_{y_1}(a)$, so the ratio of $y_1$ is $1+E_{y_1}(a)/(r\delta)$, which
is $2.14$ at $k=17$. At $k=17$ the hidden ratio is $14.45$, against $\log_2\binom{17}8=14.57$.

\subsection{Power weights}
Averaging $E_i/\Inf_i$ with weights $\Inf_i^{a+1}$ would give a good coordinate if
\begin{equation}\label{eq:power}
\sum_i\Inf_i(f)^aE_i(f)\le C\sum_i\Inf_i(f)^{a+1}
\end{equation}
held for all monotone $f$. For $a=0$ this is the monotone case of the FEI conjecture with $\sum_iE_i$ in place of $H$.

\begin{proposition}
For every $a\ne0$ there is no $C$ for which~\eqref{eq:power} holds for all monotone functions.
\end{proposition}

\begin{proof}
For $a<0$ take $F$ of Proposition~\ref{prop:middle}, with $r=\binom km$ and $\delta=2^{-k}$. The hidden coordinates
contribute $r\delta^{a+1}$ to the right side and the base coordinates $k(r\delta)^{a+1}$, a fraction $kr^a\to0$ of it.
The left side is at least $r\delta^{a}E_{z_A}$, so the quotient is at least $(E_{z_A}/\delta)/(1+kr^a)\to\infty$.
For $a>0$ let $F(x,y)=\mathrm{Maj}(y)\vee(x\wedge\bigwedge_{j\in A}y_j)$ be the function in the proof of
Proposition~\ref{prop:maxinf}, before the majority substitution, and replace each $y_j$ by a copy of the tribes
function $T(w,s)$ on fresh variables, with $s=\mathrm{round}(2^w\ln2)$. Call the result $F_w$; it is monotone. The
tribes function
has mean $\mu_w=2(1-2^{-w})^s-1$, and $\mu_w\to0$ as $w\to\infty$. The expansion of Section~2.1 applies with $x$ kept
as a block of its own, expanded in the uniform basis, and each $y_j$ expanded in the $\mu_w$-biased basis. Whether
$x\in U$ depends only on the first stage of the draw, so by homogeneity of $m$ the quantities $E_x(F_w)$ and
$\Inf_x(F_w)$ are those of $F$ computed in this mixed basis. They are continuous in $\mu_w$, and therefore converge to
$E_x(F)$ and $\Inf_x(F)=2^{-k}$. A leaf $v$ in block $j$ is pivotal for $F_w$ exactly when it is pivotal for its copy
of $T(w,s)$ and $y_j$ is pivotal for $F$, and the two events are independent. Hence
$\Inf_v(F_w)\le\Inf_v(T(w,s))=2^{1-w}(1-2^{-w})^{s-1}\le2^{1-w}$. There are $kws<kw2^w$ leaves, so they contribute at
most $kw2^w\cdot2^{(1-w)(1+a)}=2^{1+a}kw2^{-aw}$ to the right side of~\eqref{eq:power}, and this tends to $0$. The left
side is at least $\Inf_x(F_w)^aE_x(F_w)$, since every term is nonnegative. Therefore the quotient of the two sides
tends to at least $E_x(F)/\Inf_x(F)$ as $w\to\infty$, and by Proposition~\ref{prop:maxinf} this is unbounded in $k$.
\end{proof}

The proof for $a>0$ fails at $a=0$, where the leaves contribute $kws2^{1-w}(1-2^{-w})^{s-1}\sim kw\ln2$ to the right
side. Majority blocks, as in Proposition~\ref{prop:maxinf}, only reach $a>1$, because the leaf influences of
$\mathrm{Maj}_b$ sum to order $\sqrt b$ while those of $T(w,s)$ sum to order $w$. The proposition excludes uniform
averaging of the ratios ($a=-1$) and inverse-influence weights. Among power weights only $a=0$ remains, and that case
is the monotone conjecture in the form $\sum_iE_i\le C\,I$. The same construction excludes the weights
$\Inf_i\,g(\Inf_i)$ for every $g$ with $g(t)\log(1/t)\to0$ as $t\to0$: the leaf influences of $F_w$ are of order
$2^{-w}$, so the leaves contribute $O(kw\,g(2^{1-w}))$ to the right side. It says nothing about weights at or above
$\Inf_i/\log(1/\Inf_i)$, the scale of Talagrand's inequality~\cite{Tal}, which keeps $\sum_i\Inf_i/\log(1/\Inf_i)$ at
least a constant times the variance.

\begin{remark}
The spectral sample of a monotone function need not be positively correlated. For $x_1\wedge(x_2\vee x_3)$ one finds
$\Pr[1\in U]=3/4$, $\Pr[2\in U]=1/4$ and $\Pr[\{1,2\}\subseteq U]=1/8<3/16$.
\end{remark}

\section*{Acknowledgements and funding}
This work was done by multiple AI agents, Claude (Anthropic) and Codex (OpenAI), connected to Math
Brain~\cite{MathBrainRepo} and working together inside the Brain framework~\cite{BrainFramework}, a shared, persistent
memory in which the agents record claims, proofs, failed attempts and methods, retrieve one another's results,
coordinate who works on what, and audit each other's work. For more on the details behind the Brain and Math Brain, see
the Math Brain paper~\cite{MathBrainPaper}, the Math Brain repository~\cite{MathBrainRepo} and the Brain framework
repository~\cite{BrainFramework}.

The research presented in this paper was supported by the European Research Council (ERC) under the European Union's
Horizon Europe research and innovation programme (grant agreement No.~101041711), by the Simons Foundation as part of
the Collaboration on the Mathematical and Scientific Foundations of Deep Learning, by Heights Labs, by Convex Nexus
Capital, by the Israel Science Foundation (grant number 2258/19), and by the Israel Science Foundation (ISF Grant
4101/25). Computational resources were provided in part by Google Cloud through its education credits program (credit
530550103).

\section*{Declaration of competing interest}
The author holds equity in Heights Labs and in Convex Nexus Capital and has a management role at Convex Nexus Capital.
These entities are included in the funding declaration above.

\section*{Declaration of generative AI and AI-assisted technologies}
The paper comes out of Math Brain, an ongoing project on the use of AI systems for research mathematics, and its
results were developed with the Brain~\cite{BrainFramework} and Math Brain~\cite{MathBrainRepo}. In November 2026 a
Claude Fable agent (Anthropic), connected to the Brain, studied the one-coordinate constant and its composition towers
and directed Codex agents running OpenAI's GPT-6.1 on the same problem. The Codex agents ran the tower searches and the
exhaustive search over $K_{3,3}$-invariant bases, and proposed the width bound, the read-once bound and the negative
examples of Sections~4 and~5. A second Claude agent found the tower of Table~1 with the largest value and the pinned-mean
limit, and a third gave the tribes substitution that extends the power-weight result to $0<a\le1$. Claude Opus 5.5 (Anthropic) searched the literature, wrote the proofs in their present form, recomputed every
tower in Table~1 and checked every numerical claim with independent code, prepared the manuscript and audited it. These
checks are not peer review. The author directed the work and is responsible for the content of the paper and for
decisions concerning its dissemination.

\section*{Code and data availability}
\url{https://github.com/yspennstate/fei-good-coordinates}

The paper is part of the Math Brain project, whose public repository is
\url{https://github.com/yspennstate/math-brain-public}~\cite{MathBrainRepo}.

\begin{thebibliography}{99}
\bibitem{ACKSW} S.~Arunachalam, S.~Chakraborty, M.~Kouck\'y, N.~Saurabh, R.~de Wolf, Improved bounds on Fourier
entropy and min-entropy, ACM Trans. Comput. Theory 13 (2021), no.~4, Art.~22, arXiv:1809.09819.
\bibitem{BrainFramework} Y.~Shmalo, \emph{Brain Framework}, public software repository, 2026,
\url{https://github.com/yspennstate/brain-framework-public}.
\bibitem{CKLS} S.~Chakraborty, R.~Kulkarni, S.~V.~Lokam, N.~Saurabh, Upper bounds on Fourier entropy,
Electronic Colloquium on Computational Complexity TR13-052, 2013.
\bibitem{FK} E.~Friedgut, G.~Kalai, Every monotone graph property has a sharp threshold, Proc. Amer. Math. Soc. 124
(1996), 2993--3002.
\bibitem{GMP} M.~J.~Gonz\'alez, P.~MacManus, M.~C.~Pereyra, Further evidence towards the Fourier entropy-influence
conjecture, arXiv:2606.00246, 2026.
\bibitem{Han} X.~Han, A new bound for the Fourier-entropy-influence conjecture, Combinatorica 45 (2025),
arXiv:2312.08271.
\bibitem{Hod} R.~Hod, Improved lower bounds for the Fourier entropy/influence conjecture via lexicographic functions,
arXiv:1711.00762, 2017.
\bibitem{KKLMS} E.~Kelman, G.~Kindler, N.~Lifshitz, D.~Minzer, M.~Safra, Towards a proof of the
Fourier-entropy conjecture?, Geom. Funct. Anal. 30 (2020), 1097--1138.
\bibitem{KMS} N.~Keller, E.~Mossel, T.~Schlank, A note on the entropy/influence conjecture, Discrete Math. 312
(2012), 3364--3372.
\bibitem{MathBrainPaper} Y.~Shmalo, \emph{Math Brain: shared, auditable memory for multi-agent mathematical research},
October 2026, \url{https://github.com/yspennstate/math-brain-public/tree/main/papers/math-brain-2026-10-24}.
\bibitem{MathBrainRepo} Y.~Shmalo, \emph{Math Brain: public snapshot, code and tests}, 2026,
\url{https://github.com/yspennstate/math-brain-public}.
\bibitem{ODo} R.~O'Donnell, \emph{Analysis of Boolean Functions}, Cambridge University Press, 2014.
\bibitem{OT13} R.~O'Donnell, L.-Y.~Tan, A composition theorem for the Fourier entropy-influence conjecture,
ICALP 2013, arXiv:1304.1347.
\bibitem{OWZ} R.~O'Donnell, J.~Wright, Y.~Zhou, The Fourier entropy-influence conjecture for certain classes of
Boolean functions, ICALP 2011.
\bibitem{Pot26} V.~N.~Potapov, On the Fourier entropy-influence conjecture for Boolean plateaued functions,
arXiv:2609.21563, 2026.
\bibitem{Shm1} Y.~Shmalo, Read-once formulas and a question of Han on the Fourier entropy-influence conjecture,
\url{https://github.com/yspennstate/han-question-read-once-formulas}, 2026.
\bibitem{Shm2} Y.~Shmalo, The one-coordinate step in the Fourier entropy-influence conjecture,
\url{https://github.com/yspennstate/fourier-entropy-influence-one-coordinate}, 2026.
\bibitem{Tal} M.~Talagrand, On Russo's approximate zero-one law, Ann. Probab. 22 (1994), 1576--1587.
\bibitem{WWW14} A.~Wan, J.~Wright, C.~Wu, Decision trees, protocols, and the Fourier entropy-influence conjecture,
ITCS 2014, arXiv:1312.3003.
\end{thebibliography}

\end{document}
